% =============================================================================
% TECHNICAL DETAILS
% Status: Reference documentation
% =============================================================================

\subsection{Software Versions}

\begin{table}[H]
\centering
\begin{tabular}{ll}
\toprule
\textbf{Software} & \textbf{Version} \\
\midrule
FiftyOne & 1.12.0 \\
Python & 3.12 \\
R & 4.3.3 \\
PyTorch & (system version) \\
CLIP model & ViT-B/32 \\
\bottomrule
\end{tabular}
\end{table}

\subsection{Data Paths}

\begin{verbatim}
research/
+-- paper/
|   +-- main.tex
|   +-- sections/
|   |   +-- abstract.tex
|   |   +-- introduction.tex
|   |   +-- objectives.tex
|   |   +-- methods.tex
|   |   +-- results.tex
|   |   +-- discussion.tex
|   |   +-- conclusions.tex
|   |   +-- observations.tex
|   |   +-- technical.tex
|   +-- figures/
+-- data/
|   +-- raw/
|   +-- processed/
|       +-- santa_ines_data.csv
|       +-- similarity_neighbors.csv
+-- analysis/
|   +-- R/
|   |   +-- analysis.R
|   |   +-- similarity_analysis.R
|   +-- python/
+-- notes/
+-- logs/
\end{verbatim}

\subsection{Key Commands}

\subsubsection{FiftyOne Dataset Operations}

\begin{lstlisting}[language=Python]
import fiftyone as fo
import fiftyone.brain as fob

# Load dataset
dataset = fo.load_dataset("santa_ines_dron")

# Access embeddings
results = dataset.load_brain_results("clip_viz")
points = results.points  # Shape: (97, 2)

# Similarity search
view = dataset.sort_by_similarity(
    sample_id, k=10, brain_key="clip_similarity"
)

# Filter by tags
high_alt = dataset.match_tags("high-altitude")
closeups = dataset.match_tags("close-up")
\end{lstlisting}

\subsubsection{Compilation}

\begin{lstlisting}[language=bash]
# Compile LaTeX document
cd research/paper
pdflatex main.tex
pdflatex main.tex  # Run twice for TOC

# Or use latexmk for automatic compilation
latexmk -pdf main.tex
\end{lstlisting}

\subsection{FiftyOne Brain Keys}

\begin{table}[H]
\centering
\begin{tabular}{lll}
\toprule
\textbf{Brain Key} & \textbf{Type} & \textbf{Description} \\
\midrule
clip\_viz & visualization & CLIP embeddings + UMAP 2D projection \\
clip\_similarity & similarity & Similarity index for retrieval \\
\bottomrule
\end{tabular}
\end{table}

\subsection{GPS Extraction from Drone Images}

GPS coordinates were extracted from EXIF metadata embedded in DJI drone images:

\begin{lstlisting}[language=Python]
from PIL import Image
from PIL.ExifTags import TAGS, GPSTAGS

def get_gps_coords(image_path):
    """Extract GPS coordinates from EXIF data"""
    image = Image.open(image_path)
    exif_data = image._getexif()

    gps_info = exif_data.get('GPSInfo', {})

    # Convert DMS to decimal degrees
    def to_degrees(dms):
        d, m, s = dms
        return d + m/60.0 + s/3600.0

    lat = to_degrees(gps_info['GPSLatitude'])
    if gps_info['GPSLatitudeRef'] == 'S':
        lat = -lat

    lon = to_degrees(gps_info['GPSLongitude'])
    if gps_info['GPSLongitudeRef'] == 'W':
        lon = -lon

    return lat, lon
\end{lstlisting}

\subsection{Satellite Map Integration}

Google Maps Static API was used to obtain satellite imagery background:

\begin{lstlisting}[language=Python]
import requests

API_KEY = "YOUR_API_KEY"
center_lat, center_lon = -34.4746, -70.9516
zoom = 18  # High zoom for agricultural detail

url = (f"https://maps.googleapis.com/maps/api/staticmap"
       f"?center={center_lat},{center_lon}"
       f"&zoom={zoom}&size=640x640"
       f"&maptype=satellite&key={API_KEY}")

response = requests.get(url)
satellite_img = plt.imread(BytesIO(response.content))
\end{lstlisting}

\subsubsection{Coordinate Projection}

GPS coordinates were converted to image pixel coordinates using Web Mercator projection:

\begin{lstlisting}[language=Python]
import numpy as np

def gps_to_pixels(lat, lon, center_lat, center_lon,
                  zoom, img_size):
    """Convert GPS to image pixel coordinates"""
    def lat_lon_to_world(lat, lon, zoom):
        lat_rad = np.radians(lat)
        n = 2 ** zoom
        x = (lon + 180) / 360 * n * 256
        y = (1 - np.log(np.tan(lat_rad) +
             1/np.cos(lat_rad)) / np.pi) / 2 * n * 256
        return x, y

    cx, cy = lat_lon_to_world(center_lat, center_lon, zoom)
    px, py = lat_lon_to_world(lat, lon, zoom)

    x = (px - cx) + img_size / 2
    y = (py - cy) + img_size / 2
    return x, y
\end{lstlisting}

\subsection{Analysis Scripts Reference}

\begin{table}[H]
\centering
\caption{Python scripts for analysis pipeline}
\begin{tabular}{ll}
\toprule
\textbf{Script} & \textbf{Purpose} \\
\midrule
gps\_satellite\_map.py & Extract GPS from EXIF metadata \\
satellite\_map.py & Generate satellite map with GPS overlay \\
cluster\_samples.py & Select representative images per cluster \\
combined\_figure.py & Create 4-panel combined analysis figure \\
\bottomrule
\end{tabular}
\end{table}

\begin{table}[H]
\centering
\caption{R scripts for statistical analysis}
\begin{tabular}{ll}
\toprule
\textbf{Script} & \textbf{Purpose} \\
\midrule
analysis.R & Embedding visualization, cluster statistics \\
similarity\_analysis.R & Network analysis, cohesion metrics \\
\bottomrule
\end{tabular}
\end{table}
